% ============================================================
%  EXPOSÉ TEMPLATE — Single-File Version
%  Faculty of Computer Science
% ============================================================
%
%  HOW TO USE THIS TEMPLATE
%  ─────────────────────────────────────────────────────────
%  1. Replace every placeholder in [SQUARE BRACKETS] with
%     your actual information.
%  2. Guidance notes are shown in teal italics in the PDF
%     (produced by the \guidance{} command).  Delete or
%     overwrite them as you fill in your own text.
%  3. Keep the images/ folder and literature.bib next to
%     this file when compiling.
%  4. Compile with:  pdflatex → bibtex → pdflatex → pdflatex
% ============================================================

% ─────────────────────────────────────────────────────────
%  PREAMBLE
% ─────────────────────────────────────────────────────────
\documentclass[
    a4paper,
    11pt,
    DIV=11,
    bigheadings,
    idxtotoc,
    listof=totoc,
    bibtotoc,
    halfparskip,
    cleardoubleempty,
    openright,
    oneside
]{scrartcl}

% ── Language & encoding ──────────────────────────────────
\usepackage[ngerman, english]{babel}   % change first lang to set default
\usepackage[T1]{fontenc}
\usepackage[utf8]{inputenc}

% ── Layout ───────────────────────────────────────────────
\usepackage{geometry}
\geometry{a4paper, left=30mm, right=20mm}
\usepackage{setspace}
\usepackage{blindtext}
\usepackage{layout}
\setlength{\headsep}{10mm}
\setlength{\footskip}{15mm}
\setlength{\textheight}{235mm}

% ── Graphics ─────────────────────────────────────────────
\usepackage{graphicx}
\usepackage{float}
\usepackage{wrapfig}
\usepackage{subfigure}

% ── Tables ───────────────────────────────────────────────
\usepackage{longtable}
\usepackage{multirow}
\usepackage{array}
\usepackage{colortbl}
\usepackage{booktabs}

% ── Colour palette ───────────────────────────────────────
\usepackage[dvipsnames]{xcolor}
\usepackage{color}
\definecolor{haw_mag}{rgb}{0,0.112,0.47}        % THI blue – headings
\definecolor{guidancecolor}{HTML}{0D9488}        % teal – guidance text
\definecolor{barblue}{HTML}{3B82F6}
\definecolor{bargreen}{HTML}{22C55E}
\definecolor{barorange}{HTML}{F59E0B}
\definecolor{barpurple}{HTML}{8B5CF6}
\definecolor{barred}{HTML}{EF4444}
\definecolor{barteal}{HTML}{14B8A6}
\definecolor{linkgray}{HTML}{64748B}
\definecolor{gridcolor}{HTML}{E2E8F0}
\definecolor{tab_1}{RGB}{85,128,214}
\definecolor{tab_2}{RGB}{230,230,230}

% ── Heading colours ──────────────────────────────────────
\addtokomafont{section}{\color{haw_mag}\rmfamily}
\addtokomafont{subsection}{\color{haw_mag}\rmfamily}
\addtokomafont{subsubsection}{\color{haw_mag}\rmfamily}
\addtokomafont{paragraph}{\color{haw_mag}\rmfamily}
\addtokomafont{subparagraph}{\rmfamily}

% ── Hyperlinks ───────────────────────────────────────────
\PassOptionsToPackage{hyphens}{url}
\usepackage[
    pdfborder={0 0 0},
    colorlinks=true,
    linkcolor=black,
    citecolor=red
]{hyperref}

% ── Bibliography ─────────────────────────────────────────
\usepackage[numbers]{natbib}

% ── PDF inclusion ────────────────────────────────────────
\usepackage{pdfpages}

% ── Gantt chart ──────────────────────────────────────────
\usepackage{pgfgantt}

% ── Glossary ─────────────────────────────────────────────
\usepackage[toc]{glossaries}
\makeglossaries

% ── Header / footer ──────────────────────────────────────
\usepackage{scrlayer-scrpage}
\pagestyle{scrheadings}
\clearscrheadfoot
\ohead{\pagemark}
\ihead{\ifthispageodd{\leftmark}{\rightmark}}
\setheadsepline{0.4pt}
\setkomafont{pagehead}{\normalfont}
\setfootsepline{0.4pt}
\setkomafont{pagefoot}{\normalfont}
\ifoot{\footnotesize{Makram Daou}}
\ofoot{\footnotesize{Exposé}}

% ── Custom page style (used on front-matter pages) ───────
\defpagestyle{myPageStyle}{
    (0pt,0pt)
    {\hfill\pagemark}{\hfill\pagemark}{\hfill\pagemark}
    (0pt,0pt)
}{
    {(\textwidth,0.4pt)
    \footnotesize{Makram Daou}\hfill\footnotesize{Exposé}}
    {\footnotesize{Makram Daou}\hfill\footnotesize{Exposé}}
    {\footnotesize{Makram Daou}\hfill\footnotesize{Exposé}}
    (0pt,0pt)
}

% ── Guidance helper command ───────────────────────────────
%  Renders placeholder instructions in teal italics.
%  Delete \guidance{...} blocks when you fill in real text.
\newcommand{\guidance}[1]{%
    \begin{quote}
        \itshape\color{guidancecolor}\small #1
    \end{quote}
}

% ─────────────────────────────────────────────────────────
%  GLOSSARY ENTRIES  (add / remove as needed)
% ─────────────────────────────────────────────────────────
\newglossaryentry{LLM}{name={LLM},       description={Large Language Model}}
\newglossaryentry{WoZ}{name={WoZ},       description={Wizard of Oz}}
\newglossaryentry{HMI}{name={HMI},       description={Human-Machine Interface}}
\newglossaryentry{TTS}{name={TTS},       description={Text-to-Speech}}
\newglossaryentry{SDLP}{name={SDLP},     description={Standard Deviation of Lateral Position}}
\newglossaryentry{TLX}{name={NASA-TLX},  description={NASA Task Load Index}}
\newglossaryentry{UEQS}{name={UEQ-S},    description={User Experience Questionnaire, short version}}

% ─────────────────────────────────────────────────────────
%  DOCUMENT BEGIN
% ─────────────────────────────────────────────────────────
\begin{document}

\shorthandoff{"}
\widowpenalty=10000
\clubpenalty=10000
\displaywidowpenalty=10000

% ══════════════════════════════════════════════════════════
%  TITLE PAGE
% ══════════════════════════════════════════════════════════
\newgeometry{textheight=0.9\paperheight, textwidth=0.76\paperwidth,
             left=30mm, right=20mm}

\begin{titlepage}

    \begin{figure}[!h]
        \centering
        \includegraphics[width=0.4\textwidth]{images/thiRGB.jpg}
    \end{figure}

    \begin{center}
        \hrulefill
    \end{center}

    \begin{center}
        \vspace{1cm}
        \huge\textbf{Exposé}\\[2.5em]
        \normalsize
            in the Bachelor's Program User Experience Design (B.Sc.)\\
            Faculty of Computer Science\\[7em]
        \normalsize
            Thesis Title\\[1em]
        \Large\textbf{Can an LLM Match a Human Co-Driver?}\\[0.6em]
        \large Comparing In-Car Voice Assistance for German-Speaking Drivers in a Driving Simulator
    \end{center}

    \vfill

    \begin{tabular}{ll}
        First and Last Name: & \textbf{Makram Daou}        \\[1em]
        Student ID:          & 00110695                       \\[3em]
        First Supervisor:    & Prof.\,Dr.\ Ignacio Alvarez    \\[1em]
        Second Supervisor:   & [TBD]                            \\[1em]
        Submission Date:     & [TODO: submission date]                     \\
    \end{tabular}

\end{titlepage}

\restoregeometry

% ─────────────────────────────────────────────────────────
%  FRONT MATTER
% ─────────────────────────────────────────────────────────
\pagenumbering{Roman}
\thispagestyle{myPageStyle}
\tableofcontents
\cleardoublepage

% ─────────────────────────────────────────────────────────
%  MAIN CONTENT
% ─────────────────────────────────────────────────────────
\pagenumbering{arabic}
\thispagestyle{myPageStyle}

% ══════════════════════════════════════════════════════════
%  1  INTRODUCTION & MOTIVATION
% ══════════════════════════════════════════════════════════
\section{Introduction and Motivation}

\guidance{%
    \textbf{Purpose of this section:} Introduce the topic to a reader who is
    knowledgeable in the field but unfamiliar with your specific problem.
    Aim for 1--2 pages.\\[0.5em]
    \textbf{What to include:}
    \begin{itemize}
        \item \textbf{Context \& relevance} — Why does this topic matter today?
              Mention industry trends, societal challenges, or scientific gaps
              that make this research timely.
        \item \textbf{Problem hook} — Describe a concrete situation or deficiency
              that your work will address.  Use statistics or a recent event if
              available.
        \item \textbf{Scope statement} — Briefly state what the thesis will and
              will \emph{not} cover.
        \item \textbf{Structure overview} — One sentence per subsequent section
              explaining how the exposé is organised.
    \end{itemize}
    \textbf{Tip:} Write this section last, once all other sections are complete.%
}

% ══════════════════════════════════════════════════════════
%  2  PROBLEM STATEMENT & RESEARCH QUESTIONS
% ══════════════════════════════════════════════════════════
\section{Problem Statement and Research Questions}

In-car voice assistants are moving from fixed command grammars to agents built on large language models (LLMs) that interpret free-form requests and act on vehicle functions through tool calls \cite{Alvarez26Do,Kirmayr26CAR}.
Whether such agents serve drivers well remains unclear for two reasons.\glsadd{LLM}

First, existing evaluations lack a human reference point.
Benchmarks test in-car LLM agents against simulated users without any driving task, and even frontier models complete fewer than half of ambiguous requests consistently \cite{Kirmayr26CAR}.
User studies with real drivers compare LLM agents with scripted agents or with no agent at all \cite{Choe26ChatGPT,Cockram26Small}.
Neither line of work shows how far an LLM assistant falls short of a competent human assistant who has the same vehicle functions.
This difference matters beyond task success, because effortful dialogues go together with frustration and less stable lane keeping \cite{Liu26Affective}.

Second, the evidence comes almost entirely from English- or Chinese-speaking drivers.
Multilingual benchmarks show that tool-calling agents lose performance outside English, often because non-English words end up in parameters that expect fixed English values \cite{Luo26Lost,Hofman25MAPS}.
No evaluation of a tool-using in-car agent with German-speaking drivers exists.
The only in-vehicle data on German command phrasing comes from a remote pilot study with single spoken commands and no responding agent; descriptively, its German-language sessions reported the highest workload of all cells, on a sample of seven \cite{Alvarez26Do}.

The thesis therefore compares a tool-using LLM assistant with a human Wizard-of-Oz (WoZ) assistant who has the same tools and instructions, with German-speaking drivers in a driving simulator.\glsadd{WoZ}
Three research questions guide the work:

\begin{longtable}{p{0.3cm} p{0.8cm} p{12.7cm} p{0.5cm}}
    & \textbf{RQ1} & To what extent does an LLM-based in-car assistant differ from a human assistant with the same tools in task success, dialogue efficiency and error recovery when German-speaking drivers make in-cabin requests? & \\
    & & & \\
    & \textbf{RQ2} & How does the type of assistant (LLM vs.\ human) affect drivers' lateral and longitudinal vehicle control, workload, perceived usability and trust? & \\
    & & & \\
    & \textbf{RQ3} & Which request types and German-language features (compound nouns, formal and informal address, English loanwords, place names, dialect) account for failures of the LLM assistant on tasks the human assistant completes? & \\
\end{longtable}

% ══════════════════════════════════════════════════════════
%  3  OBJECTIVES
% ══════════════════════════════════════════════════════════
\section{Objectives}

The primary objective is to \emph{compare} a tool-using LLM assistant with a human assistant who has the same vehicle tools and the same instructions, on identical in-cabin tasks spoken in German by the same drivers, so that the human condition serves as an upper reference point for what the LLM assistant could achieve.
The secondary objectives support this comparison:

\begin{enumerate}
    \item \textbf{Design} a German task set of 12--16 in-cabin requests with ground-truth tool calls and completion criteria in WorlDrive, covering clear requests, requests that need a clarifying question, requests for functions the vehicle does not have, and requests that carry German-specific features (compound nouns, formal and informal address, English loanwords, place names).
    \item \textbf{Specify} the human condition as a reproducible instrument: a wizard protocol, training and a pilot that fix how the operator may use the 44 tools, when they may ask back and what they must refuse.
    \item \textbf{Conduct} a within-subject driving-simulator study with about 32 German-speaking licensed drivers on the WorlDrive platform, counterbalanced for condition order and task set.
    \item \textbf{Analyse} task success, dialogue turns, task time and response latency from the logs, lane-keeping and speed stability from the simulator telemetry, and workload, perceived usability and trust from standardised questionnaires.
    \item \textbf{Code} every LLM failure on a task the human assistant completed with a failure taxonomy (recognition, intent, tool choice, parameter, missing-capability handling, clarification) and a German-language feature scheme, with a second coder on a subset.
    \item \textbf{Contribute} the anonymised dialogue, tool-call and telemetry data of all sessions to the WorlDrive corpus, together with the task set and the wizard protocol, so that later studies can reuse the human baseline.
\end{enumerate}

% ══════════════════════════════════════════════════════════
%  4  RELATED WORK
% ══════════════════════════════════════════════════════════
\section{Related Work}

The literature was searched in the ACM Digital Library, IEEE Xplore, Google Scholar and arXiv for 2015--2026 with the keyword groups (in-vehicle OR in-car OR automotive) $\times$ (voice assistant OR conversational agent OR LLM agent) $\times$ (Wizard of Oz OR human baseline OR benchmark) $\times$ (driver distraction OR workload OR driving performance), followed by backward and forward snowballing from the three closest papers \cite{Kirmayr26CAR,Choe26ChatGPT,Large19Lessons}.
Preprints are cited only where no peer-reviewed version exists yet and are marked as such in the bibliography.
The results fall into three clusters.

\subsection{LLM-Based Assistants with Real Drivers}

Before LLMs, in-vehicle conversational agents were studied with scripted dialogue.
\citet{Braun19Service} showed in a simulator study that the assistant's personality affects how much drivers like and trust it, and \citet{Large19Lessons} distilled design guidelines for automotive conversational interfaces from a series of Wizard-of-Oz studies in which a human spoke for the system.
Recent studies replace the script with an LLM.
\citet{Choe26ChatGPT} had 40 drivers in a motion-based simulator experience no agent, a pre-scripted agent and a ChatGPT-based agent; the LLM agent produced more stable lane keeping and higher ratings of competence, affective trust and preference than the scripted one.
\citet{Cockram26Small} ran 40 drivers on a test track in a conditionally automated vehicle and found that brief LLM-initiated conversations counteract passive fatigue, with acceptance depending on the driver type.
\citet{Liu26Affective} measured 31 drivers interacting with a commercial voice assistant for navigation and showed with structural equation modelling that effortful dialogues (more turns, longer completion) go together with frustration, higher workload and poorer lateral control.
\citet{Park26Evaluating} add a methodological warning from a multi-case study of lab evaluations of generative AI: non-deterministic output, unclear system boundaries and missing logs make such evaluations hard to interpret unless model versions, prompts and every turn are recorded.

These studies agree that LLM assistants are preferred over scripted ones and that interaction effort matters for driving.
They share three limitations.
The comparator is a scripted system or no system, never a competent human, so none of them says how good an LLM assistant is in absolute terms.
Only one of them \cite{Choe26ChatGPT} lets the agent act on vehicle functions, and none scores task success against ground-truth actions.
And all of them were run in English or Chinese with native speakers of those languages.

\subsection{Agent Benchmarks, Multilingual Tool Calling and the Wizard of Oz}

Agent benchmarks measure what the user studies leave out: task success against ground truth.
$\tau$-bench \cite{Yao24Tau} introduced simulated users, domain tools and the Pass\textasciicircum{}$k$ metric for consistency across repeated trials, and found that even strong models solve fewer than half of the tasks reliably.
CAR-bench \cite{Kirmayr26CAR} transfers this design to the car: 58 interconnected vehicle tools, tasks on navigation, charging, productivity and vehicle control, and two new task types that test whether an agent asks back when a request is ambiguous and admits when a function does not exist.
Frontier reasoning models reach under 50\,\% consistent success on the disambiguation tasks and frequently act on requests they cannot fulfil.
Both benchmarks evaluate the agent against a simulated user in English and without a driving task; they state what the agent gets wrong, not what a human in the same seat would have done.

Two multilingual benchmarks explain why the language of the driver matters for exactly this kind of agent.
\citet{Luo26Lost} show for Chinese, Hindi and Igbo that tool calling fails outside English even when the model understands the intent and picks the right tool, mainly because parameter values are written in the user's language where the tool expects English values, and that no inference-time fix fully recovers English-level performance.
\citet{Hofman25MAPS} translate four agent benchmarks into eleven languages and find that performance and safety degrade in proportion to how much of the input is non-English.
WorlDrive's tools take English enumeration values and free-text names, so a German request must be mapped into them; German has not been tested in either benchmark.
The PromptDrive pilot \cite{Alvarez26Do} is the only in-vehicle data on German command phrasing: 30 remote participants in Germany and Mexico spoke 383 single commands under three priming styles, with no agent responding; the German cell held seven participants.

The Wizard of Oz is the established way to obtain a human reference for a dialogue system.
\citet{Fraser91Simulating} and \citet{Dahlback93Wizard} laid out why and how a hidden human operator can stand in for a system, and \citet{Riek12Wizard} systematically reviewed Wizard-of-Oz studies in human--robot interaction and derived reporting guidelines on wizard training, consistency and disclosure that this thesis follows.
In the automotive field the method has been used to design dialogue \cite{Large19Lessons}, not to benchmark an LLM.
The closest comparison comes from robotics: \citet{Vrins24Wizard} let 27 participants brainstorm with a robot driven once by GPT-4 and once by a wizard, found similar ratings of social intelligence, a tentative preference for the LLM, and that participants reliably identified which robot the human controlled.
That study measured perception rather than task success, and the robot had no tools to operate.

\subsection{Driver Distraction, Workload and Measurement}

Voice interaction is not free of cost for the driver.
\citet{Strayer17Smartphone} measured the cognitive workload of three commercial intelligent personal assistants while driving and found moderate to high workload that differed between assistants, which shows that the quality of the assistant, not voice interaction as such, drives the cost to the driver.
\citet{Liu26Affective} link this load to lane keeping: the standard deviation of lateral position (SDLP) and steering reversals rose with negative emotion after effortful dialogues, more than longitudinal measures did.
The instruments for the subjective measures are established: NASA-TLX for workload \cite{Hart88NASA}, the UEQ-S for perceived usability and hedonic quality, developed with a German item set \cite{Schrepp17UEQS}, and for trust either the Trust in Automation scale of \citet{Jian00Trust} or the German-validated questionnaire of \citet{Koerber19Trust}.

\subsection{The Gap}

Taken together, the user studies have real drivers but no reference point and no ground truth; the benchmarks have ground truth but simulated users, English and no driving; the multilingual benchmarks predict degradation for non-English tool calling but have not tested German or the car; and the Wizard-of-Oz literature offers the method but has not used it to benchmark an LLM assistant with tools.
No study has compared a tool-using LLM assistant with a human assistant who has the same tools and the same instructions, with German-speaking drivers in a driving task, on task success, dialogue efficiency, driving performance, workload and trust at once.
RQ1 fills the task-success and efficiency part of this gap, RQ2 the driving and experience part, and RQ3 uses the human condition to separate failures of the LLM from failures of the task.

% ══════════════════════════════════════════════════════════
%  5  METHODOLOGY
% ══════════════════════════════════════════════════════════
\section{Methodology}

The thesis uses a \textbf{Controlled Experiment}: a within-subject driving-simulator study in which every participant completes matched in-cabin tasks once with the LLM assistant and once with the human assistant.
This design fits the research questions because RQ1 and RQ2 ask for a difference between two assistants under otherwise identical conditions, and the within-subject form removes between-person variance in driving skill, language and attitude towards assistants, which is large relative to the expected effect.
RQ3 is answered by a coded analysis of the same session logs.
Four falsifiable hypotheses are tested; RQ3 is exploratory and has none.

\begin{itemize}
    \item \textbf{H1} (RQ1): The task success rate is higher with the human assistant than with the LLM assistant, and the difference is concentrated in ambiguous, missing-function and German-feature tasks rather than in clear requests.
    \item \textbf{H2} (RQ1): Completed tasks need more dialogue turns and more time with the LLM assistant than with the human assistant.
    \item \textbf{H3} (RQ2): Lateral control during a task is poorer with the LLM assistant than with the human assistant (higher SDLP relative to the pre-task baseline); no directional hypothesis is made for speed variance.
    \item \textbf{H4} (RQ2): Workload is higher and trust and perceived usability are lower after the LLM block than after the human block.
\end{itemize}

A null result on H1 for clear requests would be a finding in itself, because the benchmarks \cite{Kirmayr26CAR} predict that current models handle unambiguous requests well and break down on ambiguity and missing capabilities.

\subsection{Use Case Definition}

The object of study is the in-cabin voice assistant of a passenger car, reached through the WorlDrive research platform of the Human-Centered Intelligent Systems Lab at THI.\glsadd{HMI}
WorlDrive runs a browser-based driving simulator, records 10\,Hz driving telemetry and lets a participant hold a spoken dialogue with an assistant that can act on 44 vehicle functions (climate, seats, windows, lights, media, navigation, phone, vehicle status).
The assistant seat is filled either by an LLM agent that receives the speech transcript, replies through text-to-speech and calls the vehicle tools, or by a researcher in a Wizard-of-Oz console who hears the participant, types a reply that is spoken with the same synthetic voice, and executes the same tools.\glsadd{TTS}
Both receive the same system prompt and tool descriptions, which is what makes the human a fair reference rather than a differently briefed assistant.
Each task carries ground-truth tool calls, an initial cabin state and an expected end state, from which the platform scores success automatically.

Participants are about 32 licensed drivers who speak German as a first language or at an equivalent level, recruited at THI and in Ingolstadt; the student-heavy sample is named as a limitation.
Sessions are held in the lab with a fixed simulator setup, manual driving on a rural road with light traffic, and a push-to-talk microphone, because voice activity detection was the main source of failed utterances in the PromptDrive pilot \cite{Alvarez26Do}.
Participants are told that they will use two versions of the assistant; that one of them is a human is disclosed in the debriefing, and the ethics application covers this deception.
The LLM is one current model with tool calling and strong German performance; its name, version, system prompt, temperature and tool set are fixed in a WorlDrive agent profile before the pilot and not changed during data collection \cite{Park26Evaluating}.
[TODO: name the model once chosen with Prof.\ Alvarez and confirm who carries the API costs.]

\subsection{Independent Variables}

The single manipulated variable is the \textbf{type of assistant} with two levels, \emph{LLM agent} and \emph{human wizard}, varied within subjects.
Two matched task sets (A and B) of equal composition are built so that no participant meets the same task twice; the assignment of task set to condition and the order of conditions are both counterbalanced, giving four sequences with eight participants each at $n = 32$.
Kept constant are the route, traffic and weather, the synthetic voice, the system prompt and tool set, the task instructions on the head-up display and the experimenter.
Controlled and reported rather than manipulated are the operator identity (one or two trained wizards for all sessions, recorded per session), the response latency per condition (p50 and p95 from the platform) and the participants' driving experience and prior use of voice assistants.

\subsection{Dependent Variables}

\begin{longtable}{>{\raggedright\arraybackslash}p{3.0cm} >{\raggedright\arraybackslash}p{6.6cm} >{\raggedright\arraybackslash}p{5.2cm}}
    \toprule
    \textbf{Construct} & \textbf{Measure and unit} & \textbf{Source} \\
    \midrule
    \endfirsthead
    \toprule
    \textbf{Construct} & \textbf{Measure and unit} & \textbf{Source} \\
    \midrule
    \endhead
    Task success (RQ1, primary) & Pass/fail per task from the platform reward (correct tool calls and arguments, expected end state, policy respected); per-participant success rate per condition & WorlDrive task runs \\
    Dialogue efficiency (RQ1) & Number of user turns and time from first utterance to completion, in completed tasks only & Dialogue log \\
    Error recovery (RQ1) & Share of tasks that were completed after at least one clarification, misrecognition or wrong tool call & Dialogue and tool-call log, coded \\
    Response latency (control) & Seconds from end of utterance to start of spoken reply, p50 and p95 per condition & Platform latency report \\
    Lateral control (RQ2) & SDLP in metres and steering reversal rate during each task window, as change from the same-length pre-task baseline & 10\,Hz telemetry \\
    Longitudinal control (RQ2) & Standard deviation of speed in km/h during each task window, as change from baseline & 10\,Hz telemetry \\
    Workload (RQ2) & NASA-TLX, raw, six subscales, 0--100, after each condition block \cite{Hart88NASA} & Post-block questionnaire \\
    Perceived usability (RQ2) & UEQ-S, eight items, pragmatic and hedonic quality, $-3$ to $+3$, after each block, German original \cite{Schrepp17UEQS} & Post-block questionnaire \\
    Trust (RQ2) & Trust in Automation questionnaire, German version \cite{Koerber19Trust}, after each block; items refer to ``the assistant'' & Post-block questionnaire \\
    Failure type (RQ3) & Coded category per failed LLM task: recognition error, intent misunderstanding, wrong tool, wrong or language-mismatched parameter \cite{Luo26Lost}, acted on a missing function, missing clarification, dialogue breakdown; plus German feature present in the request & Transcripts, audio and tool calls \\
    Manipulation check & Which assistant the participant believed to be human, and confidence & Debriefing \\
    \bottomrule
\end{longtable}
\glsadd{SDLP}\glsadd{TLX}\glsadd{UEQS}

\subsection{Procedure / Research Steps}

\begin{enumerate}
    \item \textbf{Task set and wizard protocol} (weeks 1--3). Author 12--16 tasks in WorlDrive with ground truth, in two matched sets: clear single-step requests, two-step requests, requests that need a clarifying question, requests for a function the car does not have, and requests with German-specific features. Write the wizard protocol: same wiki as the LLM, may ask one clarifying question, must refuse functions that do not exist, no small talk beyond the reply. Train the wizard(s) on all tasks.
    \item \textbf{Ethics and consent} (weeks 1--4). Submit the THI ethics application with consent form, deception and debriefing plan, and the WorlDrive data-sharing terms. [TODO: confirm the route and the templates with the lab.]
    \item \textbf{Pilot} (week 5). Three to four participants run the full procedure with both conditions; check that completion fires on the intended tool calls, that the wizard's latency and style are consistent, and that the session fits 75 minutes. Adjust tasks and freeze the agent profile and the protocol.
    \item \textbf{Session} (weeks 6--10, about 32 participants): consent and demographics (10 min); practice drive without assistant (5 min); baseline drive (3 min); block 1 with assistant A, six to eight tasks spaced along the route (12 min); NASA-TLX, UEQ-S, trust (7 min); block 2 with assistant B (12 min); questionnaires (7 min); debriefing with the manipulation check and disclosure (5 min). The experimenter sits with the participant; the wizard works from a separate room and does not meet the participant before the debriefing.
    \item \textbf{Analysis} (weeks 10--13). Primary test: the per-participant task success rate is compared between conditions with a paired $t$-test, or the Wilcoxon signed-rank test if the differences are not approximately normal; as a complement, a mixed-effects logistic regression on single task outcomes with participant and task as random effects, which uses all trials and tests the interaction of condition with task type for H1. H2--H4 use paired tests on the per-participant means; the family of RQ2 outcomes is controlled with the Holm correction, and effect sizes are reported as Cohen's $d_z$ with confidence intervals \cite{Lakens13Effect}. A power analysis in G*Power \cite{Faul07GPower} for a paired comparison at $\alpha = .05$ and power $= .80$ gives $n = 34$ for a medium effect $d_z = 0.5$ and $n = 24$ for $d_z = 0.6$; 32 participants therefore detect medium-to-large effects, and the exposé states this rather than ``about 30''. For RQ3, two coders code all failed LLM tasks; agreement on a 20\,\% subset is reported as Cohen's $\kappa$.
    \item \textbf{Data release} (week 13). Pseudonymised session data, the task set and the wizard protocol are handed to the WorlDrive corpus.
\end{enumerate}

\subsection{Validity, Reliability and Ethics}

The main threats and their countermeasures are:
\emph{order and learning effects} are handled by counterbalancing condition order and task set and by the practice drive;
\emph{wizard variability} by a fixed protocol, training, one or two operators for all sessions and reporting the operator per session, following the reporting guidelines of \citet{Riek12Wizard};
\emph{latency as a confound} by measuring p50 and p95 per condition, by giving the wizard quick-reply shortcuts for routine confirmations and by reporting the difference as a limitation;
\emph{speech recognition as a confound} by noting that the LLM receives the recogniser's transcript while the wizard also hears the audio, so RQ3 separates recognition errors from reasoning and tool errors and reports both;
\emph{LLM non-determinism} by a frozen agent profile, a low temperature and full logging of every turn \cite{Park26Evaluating};
\emph{multiple outcomes} by one primary outcome and the Holm correction;
\emph{unblinding} by the manipulation check, which is analysed as a covariate.
External validity is limited by the simulator, the single model and the student-heavy sample; these are named in the limitations.

Ethics approval is obtained from THI before data collection. Participants give informed consent that covers audio recording, transcription, driving telemetry, the deception about the human assistant and the sharing of pseudonymised data in the WorlDrive corpus; they may stop at any time, including for simulator discomfort, and receive the compensation stated in the consent form. Personal data are processed under the GDPR with pseudonymous participant codes and a stated retention period. [TODO: compensation and retention period.]

% ══════════════════════════════════════════════════════════
%  6  EXPECTED CONTRIBUTIONS
% ══════════════════════════════════════════════════════════
\section{Expected Contributions and Results}

\guidance{%
    \textbf{Purpose of this section:} Communicate the anticipated value of your
    thesis.  Aim for 0.5--1 page.\\[0.5em]
    \textbf{What to include:}
    \begin{itemize}
        \item \textbf{Scientific contribution} — new knowledge, a validated
              model, a benchmark, an empirically tested framework, \ldots
        \item \textbf{Practical / applied contribution} — a prototype, tool,
              guideline, or design recommendation.
        \item \textbf{Expected findings} — your educated hypothesis about what
              you will discover (be honest about uncertainty; phrases like
              \emph{``we expect that \ldots''} or \emph{``it is hypothesised
              that \ldots''} are appropriate).
        \item \textbf{Limitations} — acknowledge foreseeable constraints
              (sample size, generalisation, time, \ldots).
    \end{itemize}
    \textbf{Tip:} Do not over-promise.  Reviewers value intellectual honesty
    over grand claims.%
}

% ══════════════════════════════════════════════════════════
%  7  WORK PLAN & SCHEDULE
% ══════════════════════════════════════════════════════════
\section{Work Plan and Schedule}

\guidance{%
    \textbf{Purpose of this section:} Demonstrate that the project is
    feasible within the allotted time.  The Gantt chart below covers a
    typical 24-week (6-month) thesis period — adjust dates and tasks to
    match your actual timeline.\\[0.5em]
    \textbf{What to include:}
    \begin{itemize}
        \item A Gantt chart (provided) showing phases, tasks, and milestones.
        \item A short paragraph (3--5 sentences) summarising the phases and
              any critical-path dependencies.
        \item Key milestones: registration, mid-point check-in, submission.
    \end{itemize}
    \textbf{Tip:} Build in buffer time before the final submission — unforeseen
    delays (data collection issues, feedback rounds) are the rule, not the
    exception.%
}

\noindent
\resizebox{\textwidth}{!}{%
\begin{ganttchart}[
    x unit=0.6cm,
    y unit title=0.55cm,
    y unit chart=0.6cm,
    canvas/.append style={draw=none},
    vgrid={*1{draw=gridcolor, thin}},
    hgrid={*1{draw=gridcolor, thin}},
    title/.style={draw=none, fill=none},
    title label font=\bfseries\footnotesize\color{black!70},
    include title in canvas=false,
    bar height=0.5,
    bar label font=\footnotesize,
    bar/.append style={draw=none, fill=barblue},
    bar incomplete/.append style={fill=barblue!30},
    group/.append style={draw=none, fill=black!70},
    group left shift=0,
    group right shift=0,
    group height=.3,
    group peaks height=0,
    group label font=\bfseries\footnotesize,
    milestone/.append style={fill=barred, draw=barred, rounded corners=1pt, xscale=0.7},
    milestone label font=\tiny\itshape,
    link/.style={-latex, thick, draw=linkgray!60},
    link bulge=4,
]{1}{24}

    \gantttitle{Thesis Schedule in Weeks}{24} \\
    \gantttitlelist{1,...,24}{1} \\

    % ── Phase 1 ────────────────────────────────────────────
    \ganttgroup{Phase 1: Topic \& Orientation}{1}{4} \\
    \ganttbar[bar/.append style={fill=barblue}]{Topic exploration \& selection}{1}{2} \\
    \ganttbar[bar/.append style={fill=barblue}]{Initial supervisor meetings}{1}{3} \\
    \ganttbar[bar/.append style={fill=barblue}]{Define research question \& scope}{2}{4} \\
    \ganttmilestone{Thesis registered}{4} \\

    % ── Phase 2 ────────────────────────────────────────────
    \ganttgroup{Phase 2: Literature Review}{4}{9} \\
    \ganttbar[bar/.append style={fill=barpurple}]{Systematic literature search}{4}{6} \\
    \ganttbar[bar/.append style={fill=barpurple}]{Reading \& annotating sources}{5}{8} \\
    \ganttbar[bar/.append style={fill=barpurple}]{Write literature review chapter}{7}{9} \\

    % ── Phase 3 ────────────────────────────────────────────
    \ganttgroup{Phase 3: Methodology \& Research}{8}{15} \\
    \ganttbar[bar/.append style={fill=barteal}]{Develop methodology / framework}{8}{10} \\
    \ganttbar[bar/.append style={fill=barteal}]{Data collection / experimentation}{10}{13} \\
    \ganttbar[bar/.append style={fill=barteal}]{Data analysis \& interpretation}{13}{15} \\
    \ganttmilestone{Midpoint check-in with supervisor}{12} \\

    % ── Phase 4 ────────────────────────────────────────────
    \ganttgroup{Phase 4: Writing}{9}{21} \\
    \ganttbar[bar/.append style={fill=bargreen}]{Write methodology chapter}{9}{12} \\
    \ganttbar[bar/.append style={fill=bargreen}]{Write results chapter}{14}{17} \\
    \ganttbar[bar/.append style={fill=bargreen}]{Write discussion \& conclusion}{17}{19} \\
    \ganttbar[bar/.append style={fill=bargreen}]{Write introduction \& abstract}{19}{21} \\

    % ── Phase 5 ────────────────────────────────────────────
    \ganttgroup{Phase 5: Revision \& Submission}{20}{24} \\
    \ganttbar[bar/.append style={fill=barorange}]{Self-revision \& proofreading}{20}{21} \\
    \ganttbar[bar/.append style={fill=barorange}]{Supervisor feedback round}{21}{22} \\
    \ganttbar[bar/.append style={fill=barorange}]{Final revisions \& formatting}{22}{23} \\
    \ganttbar[bar/.append style={fill=barorange}]{Printing \& binding}{23}{24} \\
    \ganttmilestone{Thesis submitted}{24} \\

    % ── Dependency links ───────────────────────────────────
    \ganttlink{elem3}{elem4}
    \ganttlink{elem4}{elem5}
    \ganttlink{elem6}{elem7}
    \ganttlink{elem8}{elem9}
    \ganttlink{elem9}{elem10}
    \ganttlink{elem10}{elem15}
    \ganttlink{elem15}{elem16}
    \ganttlink{elem16}{elem17}
    \ganttlink{elem18}{elem19}
    \ganttlink{elem20}{elem21}
    \ganttlink{elem21}{elem22}

\end{ganttchart}
}% end resizebox

% ─────────────────────────────────────────────────────────
%  BACK MATTER
% ─────────────────────────────────────────────────────────
\pagenumbering{Roman}
\setcounter{page}{6}

% ── Glossary ─────────────────────────────────────────────
\thispagestyle{myPageStyle}
\glssetwidest{NASA-TLX}
\setglossarystyle{alttree}
\printglossary[title=List of Abbreviations, toctitle=List of Abbreviations]
\cleardoublepage

% ── Bibliography ─────────────────────────────────────────
\thispagestyle{myPageStyle}
\bibliographystyle{dinat}
\renewcommand{\refname}{Bibliography}
\bibliography{literature}
\cleardoublepage

\shorthandon{"}

\end{document}
